\documentclass[11pt]{article}
\usepackage[margin=1in]{geometry}
\usepackage[T1]{fontenc}
\usepackage[utf8]{inputenc}
\usepackage{booktabs}
\usepackage{array}
\usepackage{hyperref}
\usepackage{xcolor}


\title{Jurassic Bitcoin: Ossified Tech Debt as an Isomorphic Mosquito for New Features\\}
\author{Patrick Dugan}
\date{April 17, 2026}

\begin{document}
\maketitle

\begin{abstract}
Bitcoin's pre-SegWit consensus history contains behaviors that are no longer
desirable transaction forms but remain informative as design motifs. We focus
on three reusable motifs: transcript multiplicity, identifier bifurcation, and
carrier camouflage. Our contribution is an observer pipeline that extracts
historical Bitcoin windows, materializes seam-family fixtures, classifies them
with a policy-envelope map, and then lifts the resulting abstractions into live
TradeLayer prototypes with BitVM-style dispute flows
\cite{linus2023bitvm,linus2026binohash,yang2021fluffy,kim2025fork}. The point
is not to revive old bugs. It is to translate fossilized Bitcoin behavior into
new usable mechanics, inspired by Robin Linus' work with BinoHash.
\end{abstract}

\section*{Motifs and Method}

We use \emph{Jurassic Bitcoin} as shorthand for the pre-SegWit historical era
in which Bitcoin still exposed many laissez-faire technical fingerprints of its
early libertarian design culture. That era is over, but some of those old
fingerprints survive as reusable protocol shapes once they are abstracted away
from the original transaction forms that carried them. In that sense, the
observer pipeline is the mosquito in amber: it preserves small but
information-rich traces of a vanished protocol ecology, then lets us recover
isomorphic mechanics without reviving the original environment itself. We can then take this DNA and produce monstrous or delightful feature modification.

We use the term \emph{Jurassic Bitcoin} for the historical period that ends, in
practical terms, before SegWit-era modularity takes over the protocol surface.
It refers to an older Bitcoin in which transaction behavior still carried many
laissez-faire, lightly curated technical fingerprints of an earlier
libertarian engineering culture: behaviors that were tolerated, under-specified, or simply left in place long enough to ossify. That period is over as a live design regime, but its fossils remain informative. The claim of this paper is that some of those neglected fingerprints are still isomorphically usable as protocol mechanics once they are abstracted away from their original,
historically contingent transaction forms. SegWit itself had its Witness Data component used in unexpected ways with Ordinals, what some would consider a predatory raptor and others would compare to a herbivorous brontosaurus, but either way the foot print has become a part of Bitcoin history: isomorphic utility is real.

If Jurassic Bitcoin supplies the fossils, the observer pipeline is the mosquito
in amber: it preserves small but information-rich traces of a vanished protocol
ecology, then lets us recover isomorphic mechanics without reviving the
original environment itself.

The central claim of this paper is that three motifs recur across both
historical Bitcoin and modern overlay design:

\begin{enumerate}
\item \textbf{Transcript multiplicity}: one broad spend family can admit
multiple effective signing transcripts.
\item \textbf{Identifier bifurcation}: externally visible identifiers can move
while the committed core digest stays fixed.
\item \textbf{Carrier camouflage}: overlay commitments can hide inside ordinary
transaction topology rather than obvious exotic carriers.
\end{enumerate}

We study those motifs with a deterministic observability toolchain and then lift them into live prototypes where BitVM-style routes, challenges,
and resolutions can be exercised without claiming that the original legacy
Bitcoin seams are themselves production-ready protocol constructions. The
resulting perspective sits downstream of BitVM's dispute-oriented execution
model and of more recent Binohash-style transaction introspection work, while
remaining methodologically closer to differential blockchain testing than to new consensus design \cite{linus2023bitvm,linus2026binohash,yang2021fluffy,kim2025fork}.

\section{Contributions}

This work makes four contributions.

\begin{enumerate}
\item It defines a reproducible \emph{historical-to-modern} workflow for turning Bitcoin fossils into measured experimental surfaces.
\item It organizes the current evidence into three reusable overlay-relevant
motifs rather than a loose list of quirks.
\item It demonstrates live prototype grafts on testnet with TradeLayer that show how these motifs can interact with BitVM-style flow graphs and dispute edges.
\item It introduces a policy-envelope mapping method for separating replay-only
Bitcoin fossils from clean modern isomorphisms.
\end{enumerate}

\section{Observability Toolchain}

The repository contributes a practical research tool rather than an alternative
node implementation. The toolchain has four layers.

\paragraph{Historical extraction.}
Mainnet windows are selected around activation boundaries and structurally
interesting transaction families, then extracted into reproducible corpora. The
current project already covers BIP16, BIP34, a 2013 payout burst, and BIP66
windows, with the pre-SegWit ``Jurassic'' tour intended to end at CSV rather
than witness-era modular Bitcoin.

\paragraph{Seam-family fixtures and replay.}
Curated fixture families capture legacy behaviors such as FindAndDelete context
splits, \texttt{SIGHASH\_SINGLE} degeneracy, and \texttt{CHECKMULTISIG}
dummy-element txid variation. These fixtures can be replayed, classified, and
reported with stable identifiers and digest-level metadata.

\paragraph{Historical carrier and sidecar benches.}
The toolchain also supports ordinary-looking 2013 mainnet carrier families,
especially payout-shaped topologies in the \texttt{279150..279320} window. That
allows overlay publication ideas to be discussed in terms of historical cover
rather than only abstract commitment formats.

\paragraph{Modern live graft harnesses.}
TradeLayer testnet is used as a controllable live environment for prototyping
procedural-token flows, oracle relays, BitVM-style caches, challenges,
resolutions, and payout edges. This gives a live isomorphism layer without
claiming that the exact historical Bitcoin transaction forms survive unchanged.
In that sense, the observer process is closer to consensus-differential systems
that use executable artifacts and cross-checking harnesses as research
instruments than to a paper-only catalogue of quirks \cite{yang2021fluffy,kim2025fork}.

\begin{figure}[t]
\centering
\fbox{\parbox[c][1.7in][c]{0.9\linewidth}{\centering
Figure placeholder.\\[0.5em]
Historical observer workflow from mainnet window extraction to fixture replay,
policy-envelope labeling, and live Litecoin graft execution.}}
\caption{Observer pipeline for the current repository. Replace this placeholder
with a workflow diagram showing extraction, replay, classification, and live
prototype grafting.}
\label{fig:observer-pipeline}
\end{figure}

\section{Three Design Motifs}

\subsection{Motif I: Transcript Multiplicity}

The first motif is that one broad spend family can admit multiple effective
signing transcripts. In the legacy seam bench, FindAndDelete subset variation
and \texttt{SIGHASH\_SINGLE} degeneracy both show that small structural changes
can alter the effective digest or collapse it without changing the broader
transaction family. The important abstraction is not the historical script quirk
itself; it is the existence of a transcript-selection manifold. Recent
introspection work such as Binohash is especially relevant here because it
re-opens the design space for deliberately engineered
transaction-observable transcript surfaces without requiring a soft fork
\cite{linus2026binohash}.

That abstraction is useful for modern overlays because many BitVM- and
DLC-adjacent systems depend on alternative off-chain transcripts, branch
packages, reveal paths, or dispute narratives. A protocol designer may want the
same higher-level state machine while preserving some freedom over transcript
selection, packaging, or attestation form.

The current repository lifts this motif into a live Litecoin relay surface where
two distinct transcript aliases are accepted against the same committed state,
and then further grafts those transcript preludes onto a challengeable
procedural-token route.

\subsection{Motif II: Identifier Bifurcation}

The second motif is that a public identifier can vary while the committed core
stays fixed. The historical dummy-element seam is the canonical example in the
current project: txids move while the legacy sighash digest remains stable. The
legacy path is policy-rejected today, but it exposes a clean separation between
\emph{semantic commitment} and \emph{public handle}.

This matters to modern overlays because practical systems often need visible
session identifiers, anchor labels, namespace handles, rendezvous tags, or
publication references even when the semantic core is fixed elsewhere. The live
Litecoin graft demonstrates the same shape with namespace-labeled relays: one
committed state, multiple public labels, and a further graft onto a live
challengeable route. Classical Bitcoin work mostly encountered this family as a
malleability problem, not as a design affordance, which makes the historical
contrast useful \cite{andrychowicz2013malleability,decker2014malleability}.

\subsection{Motif III: Carrier Camouflage}

The third motif is that overlay commitments should hide inside ordinary-looking
transaction topology rather than obviously exotic publication surfaces. The
historical 2013 carrier bench shows that strong candidate covers already existed
in ordinary payout and redistribution topologies, including high-fanout sprays,
exact-value batch redistributions, and large aggregator families.

This motif is especially relevant for oracle publication, watcher cadence, DLC
settlement sidecars, and related overlay systems. The point is not simply that
``big transactions existed.'' The point is that ordinary historical cover can
provide a model for how modern overlays think about publication shape, ambient
noise, and carrier selection. This is adjacent in spirit to transaction-shape
privacy work such as PayJoin, where the operational question is how to inhabit
ordinary-looking transaction space without leaking a distinctive surface
\cite{ghesmati2022payjoin}.

The current Bitcoin-side replay bench for synthetic oracle sidecars still needs
funded prevouts before it can be fully materialized on a live Bitcoin node.
However, the companion Litecoin prototypes already show the settlement and
dispute structure that such sidecars would want to influence.

\section{Policy-Envelope Mapping as Method}

The paper's methodological lens is a policy-envelope map rather than a fourth
coequal protocol claim. Each measured surface falls into one of three buckets:

\begin{enumerate}
\item \textbf{Replay-only fossils}: historically informative, but no longer
clean live surfaces on Bitcoin today.
\item \textbf{Ordinary historical cover}: surfaces that were already ordinary in
their native Bitcoin era.
\item \textbf{Clean modern isomorphisms}: live contemporary grafts that preserve
the same abstract motif without relying on the old transaction form.
\end{enumerate}

This classification prevents a common failure mode in ``quirk mining'' papers:
confusing a historical mechanism with a deployable modern protocol. Our
framework is explicit that many legacy Bitcoin mechanisms should remain museum
specimens, while their \emph{abstract shape} can still inform current overlay
design.

\begin{center}
\begin{tabular}{>{\raggedright\arraybackslash}p{0.21\linewidth} >{\raggedright\arraybackslash}p{0.22\linewidth} >{\raggedright\arraybackslash}p{0.22\linewidth} >{\raggedright\arraybackslash}p{0.24\linewidth}}
\toprule
Motif & Historical surface & Modern live isomorphism & Present status \\
\midrule
Transcript multiplicity & FindAndDelete context splits; \texttt{SIGHASH\_SINGLE} degeneracy & Litecoin relay aliases and fused router-dispute graft & Live as a relay and dispute prelude motif \\
Identifier bifurcation & Dummy-element txid variation & Litecoin namespace-labeled relays and fused router-dispute graft & Live as a namespace/search-surface motif \\
Carrier camouflage & 2013 payout and redistribution topologies & Litecoin procedural-token settlement/dispute graphs; Bitcoin replay bench pending funded prevouts & Historically grounded; Bitcoin live materialization still incomplete \\
\bottomrule
\end{tabular}
\end{center}

\section{BitVM Prototype Applications}

The current live prototype layer is deliberately framed as \emph{BitVM-style}
application work, not as a claim that the legacy Bitcoin quirks directly
implement BitVM \cite{linus2023bitvm}. The relevant applications are as follows.

\subsection{Motif Wrapper Versus Vanilla BitVM}

Vanilla BitVM is best understood here as a dispute-oriented execution pattern:
participants commit to a computation or state transition, and if one party
publishes an incorrect claim, the protocol exposes that claim through a
challenge path. The Jurassic Bitcoin grafts do not change that security model
and do not require changing Bitcoin Core. Instead, they add a deterministic
application wrapper around BitVM-style flows.

The wrapper has three jobs. First, transcript multiplicity lets one committed
state carry compact and full transcript packages, so a dispute graph can choose
between short relay forms, full evidence forms, or challenge-specific witness
narratives. Second, identifier bifurcation separates public handles from the
semantic core: an oracle event, Taproot Assets proof anchor, watchtower alert,
router branch, or proof package may all point back to the same committed state
without overloading one identifier. Third, carrier camouflage treats publication
as a placement problem. Instead of making every proof or sidecar look like a
special protocol announcement, the prototype attaches carrier hints such as
payout batches, rebalances, sweeps, wallet batches, or distribution shadows.

The resulting architecture is therefore not ``better BitVM consensus.'' It is a
product-facing packaging layer for BitVM-adjacent systems. A vanilla BitVM flow
provides the cache, challenge, resolution, and payout edges. The Jurassic wrapper
adds transcript aliases, namespace handles, carrier hints, procedural-token
state, and observer artifacts that make those edges easier to route, index,
audit, and integrate with TradeLayer, DLC sidecars, Lightning watchtowers,
Taproot Assets anchors, Ark-style batches, or proof-carrying execution systems.
The tradeoff is that the application layer now has more explicit watcher,
indexing, naming, and artifact-management requirements. Those requirements are
acceptable for the prototypes because they are deterministic and visible in the
flow graph, but they should not be confused with a simplification of BitVM's
underlying dispute assumptions.

\subsection{Transcript Selector and Branch Packaging}

Transcript multiplicity suggests a search bench for alternative dispute
transcripts, branch packaging, and attestation forms. In the live Litecoin
grafts, transcript aliases can be attached to the same contract reference before
the route moves into challengeable branches. Binohash is an especially relevant
comparison point because it makes transaction introspection itself into a
programmable input to BitVM-adjacent designs \cite{linus2026binohash}.

\subsection{Anchor and Session-ID Search}

Identifier bifurcation suggests search surfaces for public handles: session ids,
anchor names, rendezvous labels, or publication namespaces. The live namespace
grafts demonstrate that these handles can move while the committed state stays
fixed.

\subsection{Challengeable Carrier-Style Routing}

Carrier camouflage suggests that publication and settlement need not be thought
of as isolated \texttt{OP\_RETURN}-style acts. Instead, they can be embedded in
ordinary-looking payout and redistribution shapes. The live Litecoin prototypes
show the operational side of that claim: one route can proceed to release while
another is challenged and blocked.

\subsection{Procedural-Token Flow Graphs}

The most concrete prototype result is that the motifs can be integrated into
procedural-token flow graphs. The current fused Litecoin experiment already
supports:

\begin{itemize}
\item transcript and namespace preludes on the same contract reference,
\item a short-epoch routing phase,
\item BitVM-style cache opening,
\item challenge and resolution phases,
\item both released and blocked terminal outcomes.
\end{itemize}

This is the strongest evidence in the project that the motifs are not merely
descriptive labels but operational design patterns.

\begin{figure}[t]
\centering
\fbox{\parbox[c][1.7in][c]{0.9\linewidth}{\centering
Figure placeholder.\\[0.5em]
Fused procedural-token route with transcript and namespace preludes feeding a
BitVM-style cache, challenge, resolve, and payout split.}}
\caption{Live Litecoin flow-graph scaffold. Replace this placeholder with the
hybrid router-dispute graph used in the fused procedural-token run.}
\label{fig:fused-route}
\end{figure}

\begin{table}[t]
\centering
\begin{tabular}{>{\raggedright\arraybackslash}p{0.27\linewidth} >{\raggedright\arraybackslash}p{0.22\linewidth} >{\raggedright\arraybackslash}p{0.15\linewidth} >{\raggedright\arraybackslash}p{0.22\linewidth}}
\toprule
Prototype surface & Motif emphasis & Live status & Notes \\
\midrule
Transcript alias relay & Transcript multiplicity & Passed & Two accepted relay transcripts over one committed state \\
Identifier namespace relay & Identifier bifurcation & Passed & Stable semantic state with rotated public labels \\
Hybrid router-dispute flow & Carrier camouflage + dispute routing & Passed & Released and blocked terminal branches in one run \\
Oracle sidecar mesh & All three motifs & Passed & Alias, namespace, and payout-shaped carrier hints on one sidecar state \\
Watchtower beacon mesh & All three motifs & Passed & Compact/full alert proofs with rotating public handles \\
Statechain handoff mesh & All three motifs & Passed & Alternate handoff transcripts plus rotated checkpoint handles \\
Bitcoin oracle-sidecar replay & Carrier camouflage & Pending & Needs funded prevouts before live Bitcoin materialization \\
\bottomrule
\end{tabular}
\caption{Representative prototype outcomes. This table can be tightened into a
camera-ready results summary once final figure numbers and artifact references
are fixed.}
\label{tab:prototype-results}
\end{table}

\section{Applications Beyond BitVM}

The project began with BitVM-style dispute flows because they provide a clean
modern setting in which branch selection, relay freedom, and publication
topology can all be made visible at once. However, the same three motifs also
admit non-BitVM application meshes, and the repository now includes live
Litecoin proofs for three such cases.

\subsection{Oracle Sidecars}

The first non-BitVM application is a DLC or oracle-sidecar publication surface.
The live \emph{oracle sidecar mesh} binds one semantic sidecar state to two
distinct transcript aliases, two public namespace handles, and two
payout-shaped carrier hints. Concretely, the current LTCTEST prototype accepts
``fast'' and ``full'' attestations, rotates the public handle of the same state
between versioned sidecar labels, and associates the publication with carrier
profiles modeled on ordinary 2013 payout families. This is a direct example of
how transcript multiplicity, identifier bifurcation, and carrier camouflage can
co-exist without the application itself being framed as a BitVM dispute graph.

\subsection{Watchtower Beacons}

The second non-BitVM application is a watchtower or fraud-monitor beacon
surface. The live \emph{watchtower beacon mesh} shows that one alert state can
be expressed through compact and full proof transcripts, while the public alert
handle rotates independently of the semantic core. The same prototype also tags
the alert publication with ordinary wallet-traffic carrier hints such as
rebalance and sweep patterns. That gives a live proof that watcher systems can
use the three motifs as engineering knobs for proof packaging, search-surface
design, and publication cover.

\subsection{Statechain Handoffs}

The third non-BitVM application is a statechain-style handoff checkpoint. The
live \emph{statechain handoff mesh} treats the ownership transition as a fixed
semantic core while permitting two handoff transcripts (acknowledgement and
finalization), two rotated checkpoint handles, and two ordinary checkpoint
publication shapes. This is especially useful for the paper because it shows
that the same motifs are not confined to oracle publication or challenge
systems; they also apply to off-chain transfer and checkpoint protocols that
need observable but non-rigid coordination surfaces.

Taken together, these three meshes strengthen the paper's central claim. The
repository is no longer merely showing that old Bitcoin seams can be lifted into
BitVM-flavored demos. It now shows that the same historical abstractions can be
ported into oracle publication, watcher beacons, and statechain-style handoff
surfaces while remaining measurable, executable, and policy-envelope aware.

\section{Discussion}

The useful interpretation of these results is not that old Bitcoin bugs should
be revived. It is that historical Bitcoin reveals protocol degrees of freedom
that can be abstracted and measured. Some of those degrees of freedom are
cryptographic, such as digest aliasing or txid movement. Others are topological,
such as publication inside ordinary payout structure. The project's value lies
in making those degrees of freedom reproducible, classifiable, and portable into
modern experimental grafts.

This also clarifies the role of Litecoin in the project. Litecoin is not
presented as a substitute claim about Bitcoin consensus. It is used as a live
engineering environment where the motifs can be integrated with TradeLayer and
BitVM-style flows, demonstrating that the abstractions are fertile enough to
support working prototypes. That combination of historical corpus extraction,
replay, and live cross-environment grafting is methodologically closest to the
``observe, replay, differential-check, then generalize'' pattern seen in
consensus bug-finding work \cite{yang2021fluffy,kim2025fork}.

\section{Future Work}

The next steps fall into three categories.

\subsection{Bitcoin-side materialization}

The main remaining Bitcoin-side blocker is funded-prevout materialization for
the oracle-sidecar replay benches. Once the replay fixtures spend real inputs,
the carrier-camouflage thesis can be evaluated against live Bitcoin behavior
rather than only against historical extraction and synthetic replay structure.

\subsection{Reusable overlay templates}

The current Litecoin prototypes are already live, but they still exist as
scenario-specific harnesses. A natural next step is to encode them as reusable
flow-graph templates where relay preludes, routing buckets, and dispute edges
are declared parametrically rather than by one-off scripts.

\subsection{Binohash-style adaptation by other developers}

The broader research agenda is to encourage other Bitcoin developers to mine
these motifs in new ways. ``Binohash-style'' is the right shorthand here because
it captures the larger move of treating old script and transaction boundary
conditions as raw material for new introspection-capable constructions rather
than as dead historical accidents \cite{linus2026binohash}:

\begin{itemize}
\item alternative transcript alias surfaces for challenge systems,
\item new namespace or anchor search constructions,
\item novel uses of ordinary transaction topology as commitment cover,
\item covenant-like or contract-layer analogues of the same manifold shapes,
\item DLC and BitVM systems that make deliberate use of relay and identifier
freedom without depending on legacy policy-rejected transaction forms.
\end{itemize}

Instructionally, the message to the broader Bitcoin development community is
simple and intentionally modest. The ecosystem now includes individual
maintainers, new contributors, nonprofits, research groups, wallet teams,
infrastructure companies, and protocol startups with very different goals and
risk tolerances. This paper is not asking that community to standardize around
the specific Litecoin grafts or to rehabilitate Jurassic-era transaction forms.
It is offering a small toolkit of strange but measured abstractions, together
with application examples, and suggesting a directionally correct next step:
take these motifs, test them inside your own covenant, DLC, watchtower,
wallet-coordination, or BitVM-adjacent setting, and see whether they open new
design space without inheriting the old pathologies.

In that sense, the project's goal is not to close the book on Jurassic Bitcoin.
It is to make the fossils legible enough that other protocol engineers can adapt
their abstract affordances elsewhere.

\section*{Artifacts}

Representative repository artifacts include:

\begin{itemize}
\item \texttt{artifacts/grants/overlay\_hook\_summary.md}
\item \texttt{artifacts/grants/overlay\_carrier\_bench\_2013.md}
\item \texttt{artifacts/grants/policy\_envelope\_mapping.md}
\item \texttt{docs/grants/four\_ossification\_proposals.md}
\item \texttt{docs/litecoin-bitvm/procedural\_tokens.md}
\end{itemize}

\section*{Conclusion}

Bitcoin's old quirks are not just anecdotes. When measured carefully, they
define reproducible search surfaces. Transcript multiplicity, identifier
bifurcation, and carrier camouflage together form a compact vocabulary for
reasoning about how ossified historical behavior can inspire modern overlay
systems. The toolchain and prototypes in this repository show that such motifs
can be extracted, classified, and experimentally lifted into live BitVM-style
applications without confusing historical fossils for deployable mechanisms.

\bibliographystyle{plain}
\bibliography{jurassic_design_motifs}

\end{document}
